% !PRINTED none
% Chapter-opening page: no running head/folio visible in any of the three
% bands (consistent with an Abschnitt-opening page).

\begin{center}
\gesperrt{Sixth Section.}

\gesperrt{Groups of Linear Substitutions.}
\end{center}

\begin{center}
\rule{3cm}{0.4pt}
\end{center}

\begin{center}
§. 37.

\gesperrt{Linear Substitutions and their Composition.}
\end{center}

One of the most effective means for the formation of groups, to which at the same time many applications lead, is furnished by the linear substitutions and their composition. We have already met with such linear substitutions repeatedly, and have considered them, for example, in the second chapter of the first volume on the occasion of the law of multiplication of determinants, and then in the eleventh chapter in connection with equivalent numbers.

By a \gesperrt{linear substitution} of $n$ variables we understand a system of equations by which a system of $n$ variables $y_1, y_2, \dots y_n$ is expressed linearly in terms of another system $x_1, x_2, \dots x_n$. According to the number of the variables we distinguish \gesperrt{unary, binary, ternary, quaternary} substitutions, and in general we shall call the number of the variables the \gesperrt{dimension} of the substitution. We confine ourselves for the present to \gesperrt{homogeneous substitutions}, so that, if the coefficients are denoted by $a_i^{(\varkappa)}$, the substitution is represented by the system of equations

\weq{1}{y_\varkappa = \sum_{1,n}^{i} a_i^{(\varkappa)} x_i \qquad \varkappa = 1, 2, \dots n}

Often the variables themselves do not matter, so that such a substitution is sufficiently characterised by its coefficients $a_i^{(\varkappa)}$. One sometimes also uses
%CONT>>
